\subsubsection*{Solution}
\paragraph*{Step-by-Step Derivation}

Assume that "random from $O(d)$" means distributed according to normalized Haar measure. Let

\[
m:=d_B,\qquad p:=d_P,\qquad d=d_Bd_P=mp.
\]

Choose product bases in which the fiducial states are basis vectors. For $a\in\{1,\ldots,m\}$ and $i\in\{1,\ldots,\dim H_b\}$,

\[
V_{ai} = \sqrt p\, O_{(a,0_P),(i,0_f)}.
\]

Since the relevant row and column indices are distinct fixed subsets and Haar measure is invariant under permutations, we can suppress the fiducial indices and write

\[
V_{ai}=\sqrt p\,O_{ai}.
\]

Thus, with $M:=V^\dagger V$,

\[
M_{ij} = p\sum_{a=1}^{m}O_{ai}O_{aj}.
\]

\subparagraph*{1. Fourth moment of a Haar-orthogonal matrix}

The orthogonal Weingarten formula expresses polynomial Haar integrals as sums over pairings, with coefficients given by the inverse pairing Gram matrix; see \href{\detokenize{https://arxiv.org/abs/math-ph/0402073}}{Collins and Śniady} and \href{\detokenize{https://arxiv.org/abs/0906.4694}}{Theorem 2.1 of Banica}.

For four matrix elements there are three pairings,

\[
(12)(34),\qquad (13)(24),\qquad (14)(23).
\]

Their Gram matrix has diagonal entries $d^2$ and off-diagonal entries $d$:

\[
G=d(d-1)I+dJ,
\]

where $J$ is the $3\times3$ all-ones matrix. Therefore,

\[
G^{-1} = \frac{1}{d(d-1)}I - \frac{1}{d(d-1)(d+2)}J.
\]

Consequently, the diagonal and off-diagonal orthogonal Weingarten coefficients are

\[
\operatorname{Wg}_{\mathrm{diag}} = \frac{d+1}{d(d-1)(d+2)}, \qquad \operatorname{Wg}_{\mathrm{off}} = -\frac{1}{d(d-1)(d+2)}.
\]

Let

\[
D:=d(d-1)(d+2).
\]

For equal row indices, the fourth moment is the usual fourth moment of a uniform vector on $S^{d-1}$:

\[
\overline{O_{ai}O_{aj}O_{ak}O_{al}} = \frac{ \delta_{ij}\delta_{kl} +\delta_{ik}\delta_{jl} +\delta_{il}\delta_{jk} }{d(d+2)}.
\]

For $a\neq b$, only the row pairing $(12)(34)$ survives, giving

\[
\overline{O_{ai}O_{aj}O_{bk}O_{bl}} = \frac{ (d+1)\delta_{ij}\delta_{kl} -\delta_{ik}\delta_{jl} -\delta_{il}\delta_{jk} }{D}.
\]

\subparagraph*{2. Second moment of $M=V^\dagger V$}

Using

\[
M_{ij}=p\sum_{a=1}^{m}O_{ai}O_{aj},
\]

we have

\[
\overline{M_{ij}M_{kl}} = p^2\sum_{a,b=1}^{m} \overline{O_{ai}O_{aj}O_{bk}O_{bl}}.
\]

There are $m$ terms with $a=b$ and $m(m-1)$ terms with $a\neq b$. Hence

\[
\adjustbox{max width=\linewidth}{$\displaystyle \begin{aligned}
\overline{M_{ij}M_{kl}}
={}&
p^2m
\frac{
\delta_{ij}\delta_{kl}
+\delta_{ik}\delta_{jl}
+\delta_{il}\delta_{jk}
}{d(d+2)}
\\
&+
p^2m(m-1)
\frac{
(d+1)\delta_{ij}\delta_{kl}
-\delta_{ik}\delta_{jl}
-\delta_{il}\delta_{jk}
}{D}.
\end{aligned}$}
\]

Collecting the invariant tensors,

\[
\overline{M_{ij}M_{kl}} = \alpha\,\delta_{ij}\delta_{kl} + \beta\left( \delta_{ik}\delta_{jl} +\delta_{il}\delta_{jk} \right),
\]

where

\[
\begin{aligned}
\alpha
&=
\frac{p^2m\bigl[(d-1)+(m-1)(d+1)\bigr]}{D}
\\
&=
\frac{p^2m\bigl[m(d+1)-2\bigr]}{d(d-1)(d+2)},
\end{aligned}
\]

and

\[
\begin{aligned}
\beta
&=
\frac{p^2m\bigl[(d-1)-(m-1)\bigr]}{D}
\\
&=
\frac{p^2m(d-m)}{d(d-1)(d+2)}.
\end{aligned}
\]

Since $d=mp$, these become

\[
\alpha = \frac{ p\left[m(mp+1)-2\right] }{ (mp-1)(mp+2) },
\]

and

\[
\beta = \frac{ pm(p-1) }{ (mp-1)(mp+2) }.
\]

\subparagraph*{3. Contract with the two states}

Expand the states in the chosen real basis:

\[
|\phi\rangle=\sum_i\phi_i|i\rangle, \qquad |\psi\rangle=\sum_i\psi_i|i\rangle.
\]

Because $M$ is real and symmetric,

\[
\begin{aligned}
\left|\langle\phi|M|\psi\rangle\right|^2
&=
\sum_{i,j,k,l}
\phi_i^*\psi_j\phi_k\psi_l^*
M_{ij}M_{kl}.
\end{aligned}
\]

The three Kronecker-delta contractions are

\[
\sum_{i,j,k,l} \phi_i^*\psi_j\phi_k\psi_l^* \delta_{ij}\delta_{kl} = |\langle\phi|\psi\rangle|^2,
\]

\[
\sum_{i,j,k,l} \phi_i^*\psi_j\phi_k\psi_l^* \delta_{ik}\delta_{jl} = \langle\phi|\phi\rangle\langle\psi|\psi\rangle,
\]

and

\[
\begin{aligned}
\sum_{i,j,k,l}
\phi_i^*\psi_j\phi_k\psi_l^*
\delta_{il}\delta_{jk}
&=
\left|\sum_i\phi_i\psi_i\right|^2
\\
&=
|\langle\phi^*|\psi\rangle|^2.
\end{aligned}
\]

Therefore, for possibly unnormalized vectors,

\[
\overline{ \left|\langle\phi|V^\dagger V|\psi\rangle\right|^2 } = \alpha|\langle\phi|\psi\rangle|^2 + \beta\left( \langle\phi|\phi\rangle\langle\psi|\psi\rangle + |\langle\phi^*|\psi\rangle|^2 \right).
\]

For normalized states, $\langle\phi|\phi\rangle=\langle\psi|\psi\rangle=1$.
